\textbf{1. Kiến thức cần nhớ}
\begin{itemize}
\item \textbf{Khai phương một tích}: Với hai số $a \ge 0$ và $b \ge 0$, ta có:
$$\boxed{\sqrt{a \cdot b} = \sqrt{a} \cdot \sqrt{b}}$$
\item \textbf{Nhân các căn bậc hai}: Với $A \ge 0, B \ge 0$: $\sqrt{A} \cdot \sqrt{B} = \sqrt{A \cdot B}$.
\item \textbf{Khai phương một thương}: Với số $a \ge 0$ và số $b > 0$, ta có:
$$\boxed{\sqrt{\dfrac{a}{b}} = \dfrac{\sqrt{a}}{\sqrt{b}}}$$
\item \textbf{Chia hai căn bậc hai}: Với $A \ge 0, B > 0$: $\dfrac{\sqrt{A}}{\sqrt{B}} = \sqrt{\dfrac{A}{B}}$.
\end{itemize}

\textbf{2. Phương pháp giải các dạng}

%% DANG 9D3?8-1
\textbf{Dạng 1. Khai căn một tích, nhân các căn bậc hai}
\begin{enumerate}
\item Phân tích các số dưới dấu căn thành tích của các số chính phương ($4, 9, 16, 25, 36, 49, 64, 81, 100\dots$).
\item Áp dụng quy tắc khai phương một tích để tính từng căn bậc hai.
\item Khi nhân các căn bậc hai với nhau, nhân các biểu thức dưới dấu căn lại rồi mới khai phương.
\end{enumerate}

%% DANG 9D3?8-2
\textbf{Dạng 2. Khai căn một thương, chia các căn bậc hai}
\begin{enumerate}
\item Đổi hỗn số, số thập phân về phân số tối giản trước khi khai phương.
\item Khai phương riêng tử và mẫu: $\sqrt{\dfrac{A}{B}} = \dfrac{\sqrt{A}}{\sqrt{B}}$ (với $A \ge 0, B > 0$).
\item Với phép chia hai căn bậc hai: rút gọn phân thức dưới dấu căn $\sqrt{\dfrac{A}{B}}$ trước rồi mới tính.
\end{enumerate}

%% DANG 9D3?8-3
\textbf{Dạng 3. Tính giá trị, rút gọn biểu thức}
\begin{enumerate}
\item Kết hợp các quy tắc nhân, chia căn bậc hai cùng với các hằng đẳng thức đáng nhớ để thu gọn biểu thức.
\item Chú ý điều kiện xác định của các biến để khai căn chính xác dấu của biểu thức.
\end{enumerate}
